% Part: second-order-logic
% Chapter: sol-and-sets
% Section: power-of-continuum

\documentclass[../../../include/open-logic-section]{subfiles}

\begin{document}

\olfileid[pt]{sol}{set}{pow}

\olsection{A potência do contínuo}

\begin{explain}
Na lógica de segunda ordem, podemos quantificar sobre subconjuntos do
!!{domain}, mas não sobre conjuntos de subconjuntos do !!{domain}. Para
fazer isso diretamente, precisaríamos da lógica de \emph{terceira ordem}.
Por exemplo, se quiséssemos enunciar o teorema de Cantor, segundo o qual não
existe função !!{injective} do conjunto potência de um conjunto para o
próprio conjunto, poderíamos tentar formulá-lo como ``para todo conjunto
~$X$ e todo conjunto~$P$, se $P$ é o conjunto potência de~$X$, então não
$\cardle{P}{X}$''. Dizer que $P$ é o conjunto potência de~$X$ exigiria
formalizar que os !!{element}s de $P$ são exatamente todos os subconjuntos
de~$X$, portanto algo como $\lforall[Y][(P(Y)
  \liff Y \subseteq X)]$. O problema está em $P(Y)$: isso não é
!!a{formula} da lógica de segunda ordem, pois apenas termos podem ser
argumentos de variáveis de relação unária como~$P$.

Podemos, contudo, \emph{simular} a quantificação sobre conjuntos de
conjuntos, se o domínio for grande o bastante. A ideia é usar o fato de que
relações binárias~$R$ relacionam !!{element}s do !!{domain} com
!!{element}s do !!{domain}. Dada tal~$R$, podemos reunir todos os
!!{element}s com os quais algum~$x$ se relaciona por $R$:
$\Setabs{y \in \Domain{M}}{R(x,y)}$ é o conjunto ``codificado por''~$x$.
Reciprocamente, se $Z \subseteq \Pow{\Domain{M}}$ é alguma coleção de
subconjuntos de~$\Domain{M}$ e há ao menos tantos !!{element}s de
~$\Domain{M}$ quantos conjuntos em~$Z$, então também existe uma relação
~$R \subseteq \Domain{M}^2$ tal que todo $Y \in Z$ é codificado por
algum~$x$ por meio de~$R$.
\end{explain}

\begin{defn}
Se $R \subseteq \Domain{M}^2$, então \emph{$x$ $R$-codifica}
$\Setabs{y \in \Domain{M}}{R(x, y)}$.
\end{defn}

Se !!a{element}~$x \in \Domain{M}$ $R$-codifica um conjunto
$Z \subseteq \Domain{M}$, então um conjunto $Y \subseteq \Domain{M}$
codifica um conjunto de conjuntos, a saber, os conjuntos codificados pelos
!!{element}s de~$Y$. Assim, um conjunto $Y$ pode $R$-codificar $\Pow{X}$.
Ele o faz sse, para todo $Z \subseteq X$, algum $x \in Y$ $R$-codifica
~$Z$, e todo $x \in Y$ $R$-codifica algum~$Z \subseteq X$.

\begin{prop}
A !!{formula}
  \[
  \fn{Codes}(x, R, Z) \ident \lforall[y][(Z(y) \liff
    R(x, y))]
  \]
expressa que $s(x)$ $s(R)$-codifica $s(Z)$. A
!!{formula}
\begin{multline*}
  \fn{Pow}(Y, R, X) \ident \\
  \lforall[Z][(Z \subseteq X \lif \lexists[x][(Y(x) \land
    \fn{Codes}(x, R, Z))])] \land {} \\ \lforall[x][(Y(x) \lif
  \lforall[Z][(\fn{Codes}(x, R, Z) \lif Z \subseteq X)]]
\end{multline*}
expressa que $s(Y)$ $s(R)$-codifica o conjunto potência de~$s(X)$, isto é,
os elementos de $s(Y)$ $s(R)$-codificam exatamente os subconjuntos de
$s(X)$.
\end{prop}

\begin{explain}
Com esse artifício, podemos expressar afirmações sobre o conjunto potência
quantificando sobre os códigos dos subconjuntos, em vez de quantificar sobre
os próprios subconjuntos. Por exemplo, o teorema de Cantor agora pode ser
expresso dizendo que não existe função !!{injective} do domínio de qualquer
relação que codifique o conjunto potência de~$X$ para o próprio~$X$.
\end{explain}

\begin{prop}
A sentença
\begin{multline*}
  \lforall[X][\lforall[Y][\lforall[R][(\fn{Pow}(Y, R, X) \lif \\
      \lnot \lexists[u][(\lforall[x][\lforall[y][(\eq[u(x)][u(y)] \lif
            \eq[x][y])]] \land {}\\
        \lforall[x][(Y(x) \lif X(u(x)))])])]]]
\end{multline*}
é válida.
\end{prop}

\begin{explain}
O conjunto potência de !!a{denumerable} conjunto é !!{nonenumerable} e,
portanto, sua cardinalidade é maior que a de qualquer !!{denumerable}
conjunto (que é $\aleph_0$). O tamanho de $\Pow{\Nat}$ é chamado de
``potência do contínuo'', pois é igual ao tamanho do conjunto de pontos da
reta real,~$\Real$. Se o !!{domain} é grande o bastante para codificar o
conjunto potência de !!a{denumerable} conjunto, podemos expressar que um
conjunto tem o tamanho do contínuo dizendo que ele é equinumeroso com
qualquer conjunto~$Y$ que codifique o conjunto potência de um conjunto~$X$
de tamanho~$\aleph_0$. (Se o domínio não for grande o bastante, isto é, se
não contiver nenhum subconjunto equinumeroso com~$\Real$, tampouco poderá
haver uma relação que codifique~$\Pow{X}$.)
\end{explain}

\begin{prop}
Se $\cardle{\Real}{\Domain{M}}$, então a !!{formula}
\begin{multline*}
\fn{Cont}(Y) \ident 
\lexists[X][\lexists[R][((\fn{Aleph}_0(X) \land
      \fn{Pow}(Y, R, X)) \land \\ 
      \lforall[x][\lforall[y][((Y(x) \land {} 
      Y(y) \land \lforall[z][R(x,z) \liff R(y,z)]) \lif \eq[x][y])]])]]
\end{multline*}
expressa que $\cardeq{s(Y)}{\Real}$.
\end{prop}

\begin{proof}
  $\fn{Pow}(Y, R, X)$ expressa que $s(Y)$ $s(R)$-codifica o conjunto
  potência de~$s(X)$, que $\fn{Aleph}_0(X)$ diz ser contável. Portanto,
  $s(Y)$ é pelo menos tão grande quanto a potência do contínuo, embora possa
  ser maior (se vários elementos de $s(Y)$ codificarem o mesmo subconjunto
  de~$X$). Isso é excluído pela última conjunção, que exige que a associação
  entre elementos de $s(Y)$ e subconjuntos de $s(Z)$ por meio de $s(R)$ seja
  !!{injective}.
\end{proof}

\begin{prop}
$\cardeq{\Domain{M}}{\Real}$ sse
\begin{multline*}
  \Sat{M}{\lexists[X][\lexists[Y][\lexists[R][(\fn{Aleph}_0(X) \land \fn{Pow}(Y, R, X) \land \\
          \lexists[u][(\lforall[x][\lforall[y][(\eq[u(x)][u(y)] \lif \eq[x][y])]] \land {} \\\lforall[y][(Y(y) \lif \lexists[x][\eq[y][u(x)]])])])]]]}.
        \end{multline*}
  \end{prop}

\begin{explain}
A hipótese do contínuo é a afirmação de que o tamanho do contínuo é a
primeira cardinalidade !!{nonenumerable}, isto é, que $\Pow{\Nat}$ tem
tamanho~$\aleph_1$.
\end{explain}

\begin{prop}
A hipótese do contínuo é verdadeira sse \[\fn{CH} \ident
\lforall[X][(\fn{Aleph}_1(X) \liff \fn{Cont}(X))]\] é válida.
\end{prop}

Observe que não é verdade que $\lnot \fn{CH}$ é válida sse a hipótese do
contínuo é falsa. Em !!a{enumerable} domínio, não há subconjuntos de
tamanho~$\aleph_1$ e tampouco subconjuntos do tamanho do contínuo; por isso,
$\fn{CH}$ é sempre verdadeira em !!a{enumerable} domínio. Contudo, podemos
dar uma sentença diferente que é válida sse a hipótese do contínuo é falsa:

\begin{prop}
A hipótese do contínuo é falsa sse \[\fn{NCH} \ident
\lforall[X][(\fn{Cont}(X) \lif \lexists[Y][(Y \subseteq X \land \lnot
    \fn{Count}(Y) \land \lnot \cardeq{X}{Y})])]\] é válida.
\end{prop}

\end{document}
